\subsection{Annular ends in every residue branch}\label{subsec:fc-line-ends}

\begin{lemma}[Convergence from linear coordinates]\label{lem:fc-linear-convergence}
\lean{EnvelopeProperty.Berkovich.RelativeSpectrum.fc_tendsto_iff_linear}
\leanok
\uses{lem:fc-linear-determination}
A map along any filter converges to an actual analytic-line point if and
only if every linear-coordinate value converges to that point's value.
\end{lemma}

\begin{proof}
Analytic-line topology is the topology of polynomial evaluations. Factor
each polynomial into its leading coefficient and its finitely many linear
factors, counted with multiplicity. Convergence of the linear values and
continuity of finite products give convergence of every polynomial value.
Conversely, convergence of points implies convergence of every continuous
linear evaluation.
\end{proof}

\begin{lemma}[Convergence of a whole branch end]\label{lem:fc-branch-end-convergence}
\lean{EnvelopeProperty.Berkovich.RelativeSpectrum.fc_tendsto_branchEnd}
\leanok
\uses{lem:fc-linear-convergence}
As the value $|T-b|_y$ tends to $r$ from below, the actual point $y$ tends
to $\zeta_{b,r}$. This assertion holds uniformly over all points with the
specified coordinate values.
\end{lemma}

\begin{proof}
Fix $c$ and put $d=|b-c|$. If $d<r$, then eventually $d<|T-b|_y$;
the ultrametric equality for unequal summand values gives
$|T-c|_y=|T-b|_y\to r$. If $d\ge r$, then $|T-b|_y<r\le d$ and the
same equality gives $|T-c|_y=d$. In either case the limit is
$\max\{d,r\}=|T-c|_{\zeta_{b,r}}$. Apply linear-coordinate convergence.
\end{proof}

\begin{theorem}[Every residue branch contains a terminal annulus]\label{thm:fc-residue-end-annulus}
\lean{EnvelopeProperty.Berkovich.RelativeSpectrum.fc_exists_annulus_subset_nhds}
\leanok
\uses{lem:fc-branch-end-convergence, lem:fc-inflation-linear}
Let $r>0$ and $|b-a|\le r$. Every neighbourhood $V$ of $\zeta_{a,r}$
contains an actual annulus
$\{y:\rho<|T-b|_y<r\}$ for some $0\le\rho<r$.
In particular this holds for each standard-affinoid interior containing
that Gauss point.
\end{theorem}

\begin{proof}
The centre bound gives $\zeta_{b,r}=\zeta_{a,r}$. The preceding convergence
assertion says that the inverse-image filter of the left neighbourhoods
of $r$ under $y\mapsto|T-b|_y$ converges to this Gauss point.
Consequently $V$ belongs to that inverse-image filter. The intervals
$(\rho,r)$, $0\le\rho<r$, form a basis of left neighbourhoods at $r$,
since $r>0$. The inverse image of one such interval therefore lies in $V$,
which is the asserted annulus. For the last assertion, the interior is
open and contains the point, so it is a neighbourhood.
\end{proof}
